\documentclass[tikz,border=4pt]{standalone}
\usepackage{ctex}
\usepackage{amsmath}
\usepackage{pgfplots}
\pgfplotsset{compat=1.18}

% p8-06a: 一例速记——含参一次函数 y=(k-1)x+(2k-3) 在不同 k 值下的直线族
% 斜率 k-1，截距 2k-3：临界 k=1（斜率为零，函数退化）和 k=3/2（截距为零，过原点）
\begin{document}
\begin{tikzpicture}
\begin{axis}[
  width=9cm, height=8cm,
  axis lines=center,
  xlabel={$x$}, ylabel={$y$},
  every axis x label/.style={at={(axis cs:4.5,0)}, anchor=west},
  every axis y label/.style={at={(axis cs:0,5.5)}, anchor=south},
  xmin=-4.5, xmax=4.5,
  ymin=-5, ymax=5.5,
  xtick={-4,-3,-2,-1,1,2,3,4},
  ytick={-4,-3,-2,-1,1,2,3,4,5},
  tick style={color=black},
  ticklabel style={font=\footnotesize},
  clip=true,
]
  % k=0: y = -x - 3（斜率负、截距负 → 二、三、四象限）
  \addplot[blue!70, thick, domain=-4.5:2, samples=20] {-x - 3};
  \node[blue!70, font=\footnotesize] at (axis cs:-3.8, 1.4) {$k=0$};

  % k=1（退化）：y = 0·x + (-1) = -1（水平虚线）
  \addplot[purple!70, thick, dashed, domain=-4.5:4.5, samples=20] {-1};
  \node[purple!70, font=\footnotesize] at (axis cs:3.8, -0.45) {$k=1$（退化）};

  % k=3/2: y = 0.5x + 0（过原点 → 第一、三象限，正比例）
  \addplot[orange!90, thick, domain=-4.5:4.5, samples=20] {0.5*x};
  \node[orange!90, font=\footnotesize] at (axis cs:3.8, 2.4) {$k=\frac{3}{2}$};

  % k=2: y = x + 1（斜率正、截距正 → 一、二、三象限）
  \addplot[red!80, thick, domain=-4.5:4, samples=20] {x + 1};
  \node[red!80, font=\footnotesize] at (axis cs:3.5, 4.8) {$k=2$};

  % k=3: y = 2x + 3（斜率正、截距正，更陡）
  \addplot[red!50, thick, domain=-4:1.25, samples=20] {2*x + 3};
  \node[red!50, font=\footnotesize] at (axis cs:1.3, 5.4) {$k=3$};
\end{axis}
\end{tikzpicture}
\end{document}
